En Ethernet, la difusión jerárquica explícita frente a 1D por defecto a $P=96$ dio un cociente de tiempos 1.95\,veces y un cociente de TX agregados 2.11\,veces. Un valor mayor que 1 favorece la variante jerárquica. El máximo de la fase Bcast pasó de 5.611 a 2.474 s.


En Wi-Fi, la difusión jerárquica explícita frente a 1D por defecto a $P=96$ dio un cociente de tiempos 1.96\,veces y un cociente de TX agregados 2.13\,veces. Un valor mayor que 1 favorece la variante jerárquica. El máximo de la fase Bcast pasó de 88.904 a 40.272 s.


En Ethernet, 1D frente a SUMMA 2D, ambos con $P=64$, dio un cociente de tiempos 0.57\,veces y un cociente de TX agregados 0.53\,veces. Un valor menor que 1 favorece 1D.


En Wi-Fi, 1D frente a SUMMA 2D, ambos con $P=64$, dio un cociente de tiempos 0.59\,veces y un cociente de TX agregados 0.51\,veces. Un valor menor que 1 favorece 1D.


En la difusión aislada por Ethernet, la colectiva global frente a la jerárquica dio un cociente de tiempo 2.11\,veces y de TX agregado 4.33\,veces.


En la difusión aislada por Wi-Fi, la colectiva global frente a la jerárquica dio un cociente de tiempo 4.37\,veces y de TX agregado 4.41\,veces.


Con 16 procesos y difusión global por Ethernet, cuatro nodos frente a uno dieron un cociente 37.08\,veces.


En Ethernet, cuatro nodos con 96 procesos frente a un nodo con 24 procesos dieron un cociente $T_{96}/T_{24}$ de 9.65\,veces. Aquí un valor mayor que 1 favorece un nodo.


En Wi-Fi, cuatro nodos con 96 procesos frente a un nodo con 24 procesos dieron un cociente $T_{96}/T_{24}$ de 151.70\,veces. Aquí un valor mayor que 1 favorece un nodo.


Para $N=6144$ por Ethernet, $T_{96}/T_{24}$ fue 4.77\,veces.


Para $N=7168$ por Ethernet, $T_{96}/T_{24}$ fue 4.57\,veces.


Trapecio con $n=10^8$ y $P=24$: $T_{sim}/T_{asim}$ fue 1.70\,veces.


Trapecio con $n=10^8$ y $P=96$: $T_{sim}/T_{asim}$ fue 1.07\,veces.


Trapecio con $n=10^{11}$ y $P=24$: $T_{sim}/T_{asim}$ fue 1.05\,veces.


Trapecio con $n=10^{11}$ y $P=96$: $T_{sim}/T_{asim}$ fue 0.97\,veces.

